\chapter{安全、隔离与最小权限}

\section{原理}

操作系统的隔离不是一个单独功能，而是贯穿 CPU、内存、文件、进程、设备和网络的约束集合。没有隔离，一个进程可以读写任意内存、伪造设备请求、修改别人的文件、消耗全部 CPU 和内存。安全机制的目标不是让程序“不会出错”，而是让错误和恶意行为被限制在授权边界内。

最基础隔离来自 CPU 特权级和页表权限。用户态不能执行特权指令，不能访问内核页，不能修改页表。系统调用是受控入口，内核在入口处检查参数和权限。文件权限、用户 ID、组 ID、capability、namespace、cgroup、seccomp 和 LSM 则在更高层限制对象访问和资源消耗。

\section{实践}

安全分析先写主体、对象、操作、权限：

\begin{longtable}{p{0.18\linewidth}p{0.22\linewidth}p{0.24\linewidth}p{0.24\linewidth}}
\toprule
主体 & 对象 & 操作 & 检查 \\
\midrule
进程 & 虚拟页 & read/write/exec & 页表权限、VMA \\
进程 & 文件 & open/read/write & mode、owner、capability、mount \\
进程 & 系统调用 & exec/fork/ioctl & seccomp、LSM、资源限制 \\
容器 & CPU/内存 & 消耗资源 & cgroup quota/limit \\
\bottomrule
\end{longtable}

然后写威胁模型：攻击者控制什么输入，目标对象是什么，期望阻止什么。没有威胁模型的“安全”讨论会变成口号。

\section{算法与实现分析}

\subsection{访问控制检查}

访问控制算法接收主体、对象和操作，返回允许或拒绝。Unix mode 位是最小模型：如果主体是 owner，看 owner bits；否则如果主体属于 group，看 group bits；否则看 other bits。

\begin{lstlisting}
may_access(subject, inode, op):
  if subject.uid == inode.uid:
    bits = inode.mode.owner
  else if subject.groups contains inode.gid:
    bits = inode.mode.group
  else:
    bits = inode.mode.other
  return bits.allow(op)
\end{lstlisting}

capability 是对 root 权限的拆分。与其让 uid 0 拥有全部能力，不如把网络配置、挂载、修改时间、加载模块等能力拆开。这样服务可以只拿最小需要的能力。

\subsection{namespace 和 cgroup}

namespace 改变“看见什么”，cgroup 限制“能用多少”。PID namespace 让进程看到自己的进程树；mount namespace 改变文件系统视图；network namespace 隔离网卡、路由和端口。cgroup 限制 CPU、内存、I/O 和进程数量。

\begin{lstlisting}
container_start():
  create_namespaces(pid, mount, net, ipc, uts)
  configure_cgroup(cpu_quota, memory_limit, pids_max)
  drop_capabilities()
  apply_seccomp_filter()
  exec_init_process()
\end{lstlisting}

容器不是虚拟机。它共享宿主内核，所以内核漏洞仍可能突破容器边界。容器安全依赖 namespace、cgroup、capability、seccomp、LSM 和镜像文件系统共同工作。

\subsection{seccomp 过滤}

\code{seccomp} 允许进程限制自己未来能调用哪些系统调用。典型服务启动后，只保留 \code{read}、\code{write}、\code{epoll}、\code{futex}、\code{clock} 等必要调用，拒绝 \code{mount}、\code{ptrace}、模块加载等危险入口。

\begin{lstlisting}
seccomp_check(syscall_nr, args):
  action = bpf_program.evaluate(syscall_nr, args)
  if action == ALLOW: continue_syscall()
  if action == ERRNO: return -EPERM
  if action == KILL: terminate_process()
\end{lstlisting}

seccomp 的价值是缩小攻击面。即使攻击者控制了进程内存，也很难调用被过滤的内核入口。但过滤规则写错也会让程序在正常路径崩溃，所以必须用测试覆盖。

\section{测试}

\begin{itemize}
\tightlist
\item 用户态访问内核地址必须失败。
\item 只读映射写入必须触发权限错误。
\item 文件权限拒绝时，不创建、不截断、不写入目标文件。
\item drop capability 后，相关特权操作必须失败。
\item namespace 中不可见的对象不能通过旧路径访问。
\item cgroup 内存限制触发时，系统有明确 OOM 或失败策略。
\item seccomp 规则允许正常路径，拒绝非必要危险系统调用。
\end{itemize}

\section{源码级补充：DAC、MAC、capability、namespace、Landlock 与硬化}

\subsection{访问控制模型}

DAC 由对象拥有者决定访问，典型是 UNIX mode、uid、gid；MAC 由系统策略决定，典型是 SELinux 和 AppArmor；RBAC 按角色授权；ABAC 按属性决策。Linux 实际路径通常叠加 DAC、capability、LSM、mount flags 和 namespace。读 \filepath{fs/namei.c}、\filepath{security/security.c}、\filepath{security/selinux/} 时，要追踪同一个对象从路径解析到权限 hook 的全过程。

\begin{lstlisting}
open(path, O_WRONLY):
  resolve path to inode and mount
  check mount is writable
  check DAC mode or CAP_DAC_OVERRIDE
  call security_inode_permission()
  check file type and open flags
  create struct file
\end{lstlisting}

安全检查必须绑定实际对象。只检查字符串路径，再使用路径，是 TOCTOU 风险；解析后持有 inode、dentry 或 file 引用，再检查权限，才有稳定对象。权限失败时还要保证没有副作用：不能先 truncate 再发现权限不足。

\subsection{capability 集合}

进程有 permitted、effective、inheritable、bounding、ambient 等 capability 集合。bounding set 给进程树设置上界；ambient set 影响 exec 后保留的能力。容器安全经常要求清空不必要 capability，尤其避免 \code{CAP\_SYS\_ADMIN}、\code{CAP\_SYS\_PTRACE}、\code{CAP\_NET\_RAW}。

\begin{lstlisting}
drop_privileges():
  clear_ambient_capabilities()
  drop_bounding_set_except(required_caps)
  set_no_new_privs()
  install_seccomp_filter()
\end{lstlisting}

能力检查要放在对象所属 user namespace 下解释。\code{ns\_capable(user\_ns, CAP\_NET\_ADMIN)} 和全局 \code{capable} 不是同一个边界。容器内 root 若只在容器 user namespace 有能力，不应自动获得宿主网络或宿主挂载的管理权。

\subsection{namespace 与 cgroup v2}

\begin{longtable}{p{0.18\linewidth}p{0.56\linewidth}}
\toprule
namespace & 隔离内容 \\
\midrule
PID & 进程编号和父子视图 \\
Mount & 挂载树和路径视图 \\
Network & 网络设备、路由、端口空间 \\
IPC & System V/POSIX IPC 对象 \\
UTS & hostname/domainname \\
User & uid/gid 映射和 namespace 内 capability \\
CGroup/Time & cgroup 视图、部分时间视图 \\
\bottomrule
\end{longtable}

cgroup v2 用 \filepath{cgroup.controllers} 声明可用 controller，用 \filepath{cgroup.subtree\_control} 向子树下发 controller，用 memory、cpu、io、pids 等文件表达限制。资源限制也是安全边界：没有 pids 限制，fork bomb 会扩大到宿主；没有 memory 限制，单服务可触发全局回收和 OOM；没有 io 限制，后台写入会拖慢其他租户。

\subsection{Landlock、seccomp user notification 和 BPF LSM}

Landlock 允许非特权进程给自己加文件访问规则，是基于 LSM 的沙箱。seccomp user notification 可把某些系统调用交给监督进程决策。BPF LSM 允许加载受验证的 BPF 程序到安全 hook。这些机制的共同边界是：它们可以进一步限制自己或被管理进程，但不能凭空授予原本没有的权限。

\begin{lstlisting}
sandbox_process:
  create Landlock ruleset
  add allowed read-only paths
  restrict_self()
  install seccomp allowlist
  exec workload
\end{lstlisting}

这类沙箱必须有负面测试：禁止路径打开失败、禁止 syscall 被拒绝、规则安装后不能再获得新权限。只验证正常路径能运行，不足以证明沙箱有效。

\subsection{TPM、measured boot 和 sealed storage}

TPM 可记录启动链测量值，remote attestation 让远端验证机器启动了预期固件、bootloader、内核和策略；sealed storage 让密钥只在特定测量状态下解封。它解决的是“这台机器是不是以预期状态启动”，不是运行时所有访问控制。内核被测量启动后，仍然要依赖权限、LSM、seccomp 和更新机制维护运行时边界。

\subsection{ASLR、stack canary、CFI 与 KASLR}

ASLR 随机化栈、堆、mmap 和 PIE 基址；stack canary 检测栈溢出覆盖返回地址；CFI 限制间接控制流；KASLR 随机化内核基址。它们是缓解措施，不是权限模型替代品。

\begin{longtable}{p{0.2\linewidth}p{0.36\linewidth}p{0.32\linewidth}}
\toprule
机制 & 防什么 & 不能防什么 \\
\midrule
ASLR/KASLR & 降低地址可预测性 & 信息泄漏后效果下降 \\
stack canary & 简单栈返回地址覆盖 & 任意写、逻辑漏洞 \\
CFI & 非法间接跳转 & 合法路径上的权限绕过 \\
W\^X/NX & 数据页执行 & JIT 配置错误、ROP \\
\bottomrule
\end{longtable}

安全测试要同时证明正常路径可用、禁止路径失败、失败有日志、敏感数据不泄露。只打开硬化选项但没有负面测试，不能说明边界有效。

\subsection{open 权限检查不能只看路径字符串}

路径是输入，inode 和 file 才是稳定对象。安全代码若先检查字符串路径，再在另一步打开文件，中间目录项可能被替换，形成 TOCTOU。更可靠的做法是让内核路径解析在同一次操作里完成权限、mount、LSM 和 open flag 检查，或者使用目录 fd、\code{openat2} 约束和已经打开的 fd 作为稳定能力。

\begin{lstlisting}
unsafe pattern:
  if path starts with "/safe":
    open(path)

safer pattern:
  dirfd = open("/safe", O_PATH | O_DIRECTORY)
  openat2(dirfd, relative_path,
          RESOLVE_BENEATH | RESOLVE_NO_SYMLINKS)
\end{lstlisting}

这里的关键不是某个 API 名字，而是对象绑定：检查和使用必须落在同一个解析结果上。符号链接、bind mount、rename、并发 unlink、mount namespace 都会改变“同一个字符串”指向的对象。

\subsection{exec 时 capability 和 no\_new\_privs 改变边界}

权限不是进程启动后固定不变。exec 会重新计算凭据、环境、dumpable 状态、secureexec、文件 capability 和 setuid 影响。若进程设置 \code{no\_new\_privs}，后续 exec 不能获得比当前更多的权限；seccomp 过滤器通常要求先设置这个标志，防止低权限进程安装过滤器后再通过 exec 提权。

\begin{lstlisting}
exec credential sketch:
  read current cred
  inspect executable mode and file capabilities
  if no_new_privs:
    suppress privilege gain
  compute new permitted/effective/ambient sets
  apply secureexec rules if privilege boundary changed
  commit new cred at exec commit point
\end{lstlisting}

ambient capability 容易被误用。它能跨 exec 传递给非特权可执行文件，所以启动服务前应清理不需要的 ambient set，并收紧 bounding set。否则子进程可能在没预期的路径上保留网络、ptrace 或管理能力。

\subsection{seccomp 规则要按参数和返回动作设计}

只按 syscall 号 allow/deny 往往太粗。例如 \code{ioctl}、\code{prctl}、\code{clone}、\code{socket} 的危险性高度依赖参数。高质量 seccomp profile 至少要区分常用参数组合，并为拒绝路径选择明确动作：返回 errno、kill 进程、记录日志或交给 user notification。

\begin{lstlisting}
seccomp policy sketch:
  allow read/write/close/futex/clock_gettime
  allow mmap only without PROT_EXEC|PROT_WRITE together
  allow socket only AF_UNIX if service does not need network
  reject mount/ptrace/bpf/module loading
  return EPERM for expected probes
  kill process for impossible dangerous calls
\end{lstlisting}

规则还要随库和运行时变化测试。动态链接器、线程库、DNS、日志库都可能调用你没想到的系统调用。部署前应在真实启动路径、热路径、错误路径和 reload 路径下跑 profile，而不是只跑主函数。

\subsection{Landlock 的边界：只能收紧，不能授权}

Landlock 适合让普通进程给自己加文件系统沙箱。它的 ruleset 声明要处理哪些访问类型，再把路径 fd 加入 allow list，最后 \code{restrict\_self}。一旦限制生效，进程不能用 Landlock 获得原来没有的权限，只能在原权限基础上进一步收紧。

\begin{lstlisting}
landlock flow:
  ruleset = create_ruleset(handled_access_fs)
  dirfd = open(allowed_dir, O_PATH | O_DIRECTORY)
  add_path_beneath_rule(ruleset, dirfd, allowed_access)
  prctl(PR_SET_NO_NEW_PRIVS, 1)
  landlock_restrict_self(ruleset)
  exec or run workload
\end{lstlisting}

测试要覆盖允许目录、禁止目录、符号链接、rename、临时文件和继承。若只测试“程序还能启动”，无法说明沙箱真的挡住了写入、删除、执行或路径逃逸。

\subsection{BPF LSM 和传统 LSM 的关系}

LSM hook 把安全检查插入到 inode、file、task、socket、bpf 等对象操作中。传统 LSM 如 SELinux、AppArmor 使用预先加载的策略；BPF LSM 允许受验证的 BPF 程序挂到安全 hook 上，适合做可观测、渐进式或环境相关的策略。

\begin{lstlisting}
security_inode_permission(inode, mask):
  for each registered LSM:
    ret = lsm->inode_permission(inode, mask)
    if ret != 0:
      return ret
  return 0
\end{lstlisting}

BPF LSM 仍受 verifier、权限和 hook 上下文限制。程序不能任意睡眠，不能随便访问内核内存，返回值语义也要遵守 hook 约定。它是内核安全框架的一部分，不是绕过 DAC/capability 的后门。

\subsection{容器逃逸常见入口}

容器共享宿主内核，所以边界经常坏在“多给了一个能力或设备”。常见风险包括挂载宿主敏感路径、保留 \code{CAP\_SYS\_ADMIN}、开放容器引擎 socket、允许特权设备、未限制 \code{/proc} 和 \code{/sys}、seccomp profile 过宽、cgroup pids/memory 未限。

\begin{longtable}{p{0.25\linewidth}p{0.36\linewidth}p{0.27\linewidth}}
\toprule
风险 & 为什么危险 & 收紧方向 \\
\midrule
\code{CAP\_SYS\_ADMIN} & 覆盖大量管理操作 & 默认删除，只按具体需要添加 \\
宿主容器引擎 socket & 等价于控制宿主容器引擎 & 不挂载，改用受限代理 \\
特权设备节点 & 可能直接访问内核或硬件接口 & devices cgroup 和最小设备集 \\
宽松 \code{/proc} & 泄露进程、内核和 namespace 信息 & hidepid、只读挂载、裁剪视图 \\
无 pids/memory 限制 & fork bomb 或内存压力拖垮宿主 & cgroup v2 限额 \\
\bottomrule
\end{longtable}

安全隔离要按“最小必要对象”配置，而不是按“程序跑起来需要什么就全开”。每放开一个 capability、设备、挂载或 syscall，都要写清楚它对应的业务需求和负面测试。
